\chapter{Jupyter Kernel}

\section{What the Jupyter kernel is}

\hayashi{} ships with a native Jupyter kernel called \texttt{hay-kernel}. It
lets you run \hayashi{} code inside Jupyter notebooks, JupyterLab, or any frontend
compatible with the Jupyter protocol.

Unlike the REPL, the kernel persists state across cells: variables, functions,
and data defined in one cell remain available in subsequent cells of the same
notebook.

\section{Building and installing}

The kernel is part of the \hayashi{} repository and is built with the
\texttt{native} feature flag:

\begin{lstlisting}[style=inline, language=sh]
cd hayashi
cargo build --release --bin hay-kernel --features native
\end{lstlisting}

Then install the kernel spec:

\begin{lstlisting}[style=inline, language=sh]
./target/release/hay-kernel --install
\end{lstlisting}

This registers the kernel under
\texttt{\$HOME/.local/share/jupyter/kernels/hayashi/} (or the equivalent
platform-specific user data directory).

Verify that the kernel is registered:

\begin{lstlisting}[style=inline, language=sh]
jupyter kernelspec list
\end{lstlisting}

The output should contain a line like:

\begin{lstlisting}[style=inline]
hayashi    ~/.local/share/jupyter/kernels/hayashi
\end{lstlisting}

\section{Using the kernel}

\subsection{JupyterLab}

Start JupyterLab:

\begin{lstlisting}[style=inline, language=sh]
jupyter lab
\end{lstlisting}

Create a new notebook and select the \textbf{Hayashi} kernel from the launcher or
the kernel menu.

\subsection{Jupyter Notebook}

Start Jupyter Notebook:

\begin{lstlisting}[style=inline, language=sh]
jupyter notebook
\end{lstlisting}

Create a new notebook and choose \textbf{Hayashi} from the kernel list.

\subsection{Running cells}

Code cells accept \hayashi{} syntax directly:

\begin{lstlisting}[caption={Example cell in Jupyter}]
let x = 1 + 2
print(x)
\end{lstlisting}

The output \texttt{3} appears below the cell. Variables persist across cells:

\begin{lstlisting}[caption={Later cell in the same session}]
let y = x * 10
print(y)
\end{lstlisting}

\subsection{Jupyter Console}

If it is installed, the console can also be used:

\begin{lstlisting}[style=inline, language=sh]
jupyter console --kernel=hayashi
\end{lstlisting}

\section{Uninstalling}

Remove the kernel spec directory:

\begin{lstlisting}[style=inline, language=sh]
rm -rf ~/.local/share/jupyter/kernels/hayashi
\end{lstlisting}

\section{Limitations in hosted environments}

Google Colab, Kaggle Notebooks, and similar hosted environments do not easily
allow registering custom Jupyter kernels, and the runtime is ephemeral. In those
cases the practical workaround is to upload a pre-built \texttt{hay} binary and
call it from Python cells, optionally through a cell magic such as
\texttt{\%\%hay}.
